\documentclass[11pt]{article}

% ------------------------------------------------
% Packages
% ------------------------------------------------
\usepackage[margin=1in]{geometry}
\usepackage{amsmath}
\usepackage{xcolor}
\usepackage{setspace}
\usepackage{tcolorbox}
\usepackage{enumitem}
\usepackage{hyperref}
\usepackage{amssymb}

\onehalfspacing

% ------------------------------------------------
% Formatting environments
% ------------------------------------------------

% Referee comment box
\newtcolorbox{commentbox}{
colback=gray!10,
colframe=gray!40,
boxrule=0.5pt,
arc=2pt,
left=6pt,
right=6pt,
top=6pt,
bottom=6pt
}

% Author response
\newenvironment{response}
{\begin{quote}\color{black}}
{\end{quote}}

% Quoted manuscript change
\newtcolorbox{revisionbox}{
colback=white,
colframe=black,
boxrule=0.5pt,
arc=2pt,
left=6pt,
right=6pt,
top=6pt,
bottom=6pt
}

\newenvironment{subcaption2}
{\strut
\vspace{-5pt}
\begin{minipage}[b]{0.95\textwidth}
  \hspace*{-\parindent}
  \footnotesize}
 {\end{minipage}}

\usepackage[backend = biber,
            style=authoryear-icomp,
            citestyle=authoryear-comp,
            url=false,
            isbn=false,
            doi=true,
            % doi=false,
            maxcitenames=2,
            uniquelist=false,
            date=year,
            uniquename=false,
            sorting=ynt,
            sortcites]{biblatex}

\addbibresource{5_Literature/Literature.bib} 

% ------------------------------------------------
% Title
% ------------------------------------------------

\title{Response to Editor \\
\textbf{JEEM-D-25-01276R1} \\ 
The heterogeneous effects of climate policy on households: Evidence from 88 countries\footnote{Initially submitted as "Compensation design for carbon pricing with horizontal heterogeneity: Evidence from 88 countries".}
}

\author{Leonard Missbach, Jan Christoph Steckel}

\date{\today}

\begin{document}

\maketitle

\noindent Dear Erica Myers,

\noindent We would like to thank you for your careful handling of the review process and your thoughtful editorial judgment. We specifically appreciate your constructive editorial decisions and suggestions, which have indeed helped us improve our work. 

\noindent Below, we reply to your comments in greater detail. 

\noindent Thank you again,

\noindent Leonard Missbach and Jan Christoph Steckel

\newpage

% =================================================
% EDITOR
% =================================================

\section*{Response to the Editor}

\subsection*{General comment}\label{comment_E_0}

\begin{commentbox}
    Thank you for the careful and thoughtful revisions. I think the paper is substantially stronger in its revised form. In particular, I found the new Theil decomposition analysis to be a meaningful improvement that helps sharpen the distinction between within- and between-income-group heterogeneity and addresses an important concern raised in the first-round reviews. I also appreciate the efforts to streamline parts of the manuscript and to clarify the interpretation and limitations of the empirical approach throughout.  

    
\end{commentbox}

\begin{response}
    We really appreciate your suggestions, which have helped us improve the manuscript.
\end{response}

\begin{commentbox}
    Given how thoroughly you addressed both the reviewers’ and my concerns, I have decided not to return the paper to the referees in order to streamline the review process. The main issue that still stands out to me concerns the clustering discussion. I appreciate that you have already toned down the interpretation of the clustering exercise, but I think the paper would benefit from going further in this direction. Both reviewers expressed skepticism about this part of the paper in the first round, and while I agree the revisions improved the presentation, I still do not think the clustering exercise is as compelling or as central as the earlier supervised ML analysis.

     In particular, I encourage you to substantially streamline the clustering discussion in the main text. My suggestion would be to frame the clustering exercise primarily as a descriptive synthesis of broad cross-country patterns in predictor importance, rather than as evidence of sharply delineated country “types.” I would then move most of the cluster-by-cluster interpretation and country-specific discussion to the appendix so that the main text remains more tightly focused on the core supervised ML results. Right now, for a descriptive exercise, this section occupies a substantial amount of space, and the discussion can be difficult to follow because many references are made to clusters that are only defined in the appendix. Along the same lines, I would encourage you to remove or substantially reduce references to the clustering exercise in the abstract, introduction, and conclusion.

     Something along the lines of the following would likely suffice (no need to follow this wording exactly, but to give you a sense of what I'm envisioning):

    “To summarize broad similarities in the importance of predictors across countries, we perform a descriptive clustering exercise using adjusted feature importance measures. The goal is not to identify sharply delineated country types, but rather to provide a compact summary of cross-country patterns in the predictors of carbon intensity.”
\end{commentbox}

\begin{commentbox}
   

    Then perhaps a brief paragraph summarizing the broad patterns:

    “The clustering exercise suggests substantial cross-country heterogeneity in the observable predictors of carbon intensity. In some groups of countries, expenditures and spatial characteristics explain relatively larger shares of variation, while in others, variables related to transportation, electricity access, or household appliances appear more important. However, predictive performance and feature availability vary substantially across countries, so these clusters should be interpreted cautiously and primarily as descriptive summaries.”

    I would then point readers to the appendix for detailed cluster memberships and robustness analyses.   
\end{commentbox}

\begin{response}
    We acknowledge that the clustering exercise can be difficult to follow. We respond to your suggestions by framing the clustering exercise constantly as \textit{descriptive}, \textit{stylized} or \textit{illustrative}. We have shortened the corresponding sections in both the Methods and Results section by shifting text to the Appendix. 

    The corresponding paragraph in the method section now reads as follows (line 469 and following):   
\end{response}

\begin{revisionbox}
    \paragraph{Identification of descriptive country clusters} Country-level analyses can be useful to identify country-specific household characteristics that are associated with higher carbon intensity of consumption. To describe similarities and differences in the importance of characteristics across many countries, we seek to identify illustrative clusters of countries. We do not aim at identifying sharply delineated country types, but rather at providing a stylized summary of cross-country patterns concerning the most important features for predicting carbon intensity of consumption.

    We adjust for the importance of individual features using country-level goodness of fit (R\textsuperscript{2}) to account for differences in available features across countries (see also Figure A.1.3). Based on the (adjusted) feature importance measure for each country, we use the k-means algorithm \autocite{MacQueen.1967} for clustering on normalized feature values. Using average silhouette widths \autocite{Rousseeuw.1987}, we identify $k=10$ as the optimal number of clusters. Within the same cluster, individual features are similarly important in predicting households' carbon intensity. We use such descriptive country clusters primarily for visualization of feature importance. More information on the clustering exercise and supplementary robustness checks can be found in Appendix A.4. 
\end{revisionbox}

\begin{response}
    The corresponding paragraph in the results section now reads as follows (line 713 and following):
\end{response}

\begin{revisionbox}
    \paragraph{Descriptive country clusters}
    For illustrative purposes, we can compare countries with respect to features that are important in predicting differences in carbon intensity. Based on the adjusted importance of all features and the vertical distribution coefficient of countries, we identify ten distinct stylized clusters of countries. Within clusters, countries are more similar to each other than to countries in other clusters. 
    
    This illustrative clustering exercise suggests substantial heterogeneity across countries in contribution of observable features to the prediction of carbon intensity in consumption (Figure A.10 and Appendix A.4). In some clusters of countries, household expenditures and spatial characteristics explain comparatively larger shares of variation in our outcome variable. In others, variables describing vehicle ownership, access to electricity, or information about primary fuel use appear to be more relevant. Implicitly, these stylized clusters indicate heterogeneous patterns of consumption and energy use. One implication is that it may be inaccurate to infer the distributional effects of climate policy in one country from the experience of other countries.

    Importantly, the model performance and feature availability can vary substantially across countries, which confines our clustering approach to primarily descriptive purposes.
\end{revisionbox}

\begin{response}
    Following your suggestion, we have removed references to the clusters in the discussion. 

    Our impression is that discussing alternative compensation options works best with application to specific countries and our country-level analyses. We thus retain such references, but omit pointing to stylized clusters as suggested. 
    
    We also propose to keep the clusters for visual representation purposes (in Figure 3) -- otherwise, comparing country-level results could be difficult for the readers.
    
    As suggested, we have made sure to remove or reduce references to the clustering exercise in abstract, introduction and conclusion.
\end{response}

\begin{commentbox}
    As a separate clarification point, I think Reviewer 2’s suggestion regarding pooled estimation with country indicators is conceptually reasonable, since flexible tree-based ML methods can accommodate country-specific nonlinear relationships through interactions. At the same time, I appreciate the practical challenges posed by heterogeneous survey instruments, differing feature availability across countries, and the national focus of the policy exercise. I do not think additional estimation is necessary at this stage, but I wanted to clarify that the rationale for separate country estimation primarily reflects concerns about comparability, interpretability, and feature heterogeneity, rather than an inability of pooled ML approaches to accommodate heterogeneous relationships.
\end{commentbox}

\begin{response}
    Yes, it is correct that boosted regression trees can accommodate country-specific non-linear relationships. We also agree that this approach could yield insights that are informative following the arguments put forward in our manuscript. In response, we have decided to add a supplementary analysis including all countries of the European Union (plus Norway and United Kingdom). The resulting analysis informs about which features predict household carbon intensity of consumption across all countries. We have added the corresponding analysis to the Appendix (see Figure X) and refer to this analysis in the text in footnote 39 in line 845 as follows:
\end{response}

\begin{revisionbox}
    Our analysis provides a foundation for more comprehensive analyses using more nuanced data\footnote{For example, fitting a model on observations from multiple countries including country as a feature can have its merits, if regulation affects the population across many countries. For illustrative purposes, we have fitted such a model to countries of the EU (see Figure \ref{}). Results show that heterogeneity could be reduced most effectively by transfers differentiated by country.}.
\end{revisionbox}

\begin{response}
    We address the issue of missing features by coding missing information explicitly. Interactions with country-level features should address the issue of missing information. The resulting model performance is 0.2\% (R\textsuperscript{2}). The most important features are household country (43\%), household expenditures (19\%) and main heating fuel (8\%).   
\end{response}

\clearpage

\end{document}